% =====================================================================
%  ALANYA — Dossier de conception
%  Étape 1 : socle de compte unifié
%
%  Compilation :  latexmk -pdf socle-compte.tex
%              |  pdflatex socle-compte.tex  (deux passes)
%              |  tectonic socle-compte.tex
% =====================================================================
\documentclass[11pt,a4paper,oneside]{report}

\usepackage{alanya-conception}
\setAlanyaDocTitre{Socle de compte}

\hypersetup{
  colorlinks=true,
  linkcolor=AlanyaIndigoStrong,
  urlcolor=AlanyaIndigoStrong,
  pdftitle={Alanya - Socle de compte unifie},
  pdfauthor={Equipe Alanya},
  pdfsubject={Conception - socle de compte unifie (types de compte et verification)}
}

\begin{document}

% =====================================================================
%  PAGE DE GARDE
% =====================================================================
\begin{titlepage}
\thispagestyle{empty}
\begin{tikzpicture}[remember picture, overlay]
  \fill[AlanyaIndigo]
    (current page.north west) rectangle ([yshift=-6.4cm]current page.north east);
  \fill[AlanyaIndigoDark]
    ([yshift=-6.4cm]current page.north west) rectangle ([yshift=-6.0cm]current page.north east);
\end{tikzpicture}

\vspace*{1.1cm}
{\color{white}\sffamily
\begin{center}
  {\small\bfseries DOSSIER DE CONCEPTION}\\[6pt]
  {\footnotesize\color{AlanyaContainer} Document de travail --- à valider avant développement}
\end{center}}

\vspace{3.4cm}

\begin{center}\sffamily
  {\small\bfseries ÉTAPE 1 --- FONDATION}\\[10pt]
  {\color{AlanyaLine}\rule{0.6\linewidth}{0.6pt}}\\[20pt]
  {\fontsize{32}{36}\selectfont\bfseries\color{AlanyaIndigoDark} Socle de compte unifié}\\[10pt]
  {\large\color{AlanyaGray} Types de compte, vérification et badges}\\[6pt]
  {\normalsize\itshape\color{AlanyaGray}
   Fondation commune aux comptes business, à la certification\\ et au compte officiel}\\[22pt]
  {\color{AlanyaLine}\rule{0.6\linewidth}{0.6pt}}
\end{center}

\vspace{1.2cm}

\begin{center}
  \badgePastilleCoche[2.6]\hspace{1.1cm}
  \badgePanier[2.6]\hspace{1.1cm}
  \badgePastillePanier[2.6]\hspace{1.1cm}
  \badgeMarque[2.6]
\end{center}

\vfill

\begin{center}\sffamily
  \begin{tabular}{@{}r@{\quad}l@{}}
    \thd{Périmètre}   & Talky (Flutter) \textbullet{} Alanya-Backend \textbullet{} Alanya Admin\\[3pt]
    \thd{Statut}      & Conception --- aucune ligne de code écrite\\[3pt]
    \thd{Précède}     & Listes de contacts, compte officiel, comptes business, certification\\[3pt]
    \thd{Date}        & \today\\
  \end{tabular}
\end{center}

\vspace{0.8cm}
\begin{center}\color{AlanyaGray}\small\sffamily
  Document confidentiel --- diffusion restreinte
\end{center}
\end{titlepage}

% =====================================================================
%  AVANT-PROPOS
% =====================================================================
\chapter*{Avant-propos}
\addcontentsline{toc}{chapter}{Avant-propos}
\markboth{Avant-propos}{}
\thispagestyle{plain}

Quatre fonctionnalités ont été demandées pour Alanya : les comptes business, la
certification des comptes, un compte officiel au nom de l'application, et les listes de
contacts préférés.

Trois d'entre elles --- business, certification, compte officiel --- ne sont pas trois
systèmes différents. Ce sont trois réponses à la même question : \emph{qui est ce compte ?}
Les concevoir séparément produirait trois systèmes de badges concurrents, trois écrans de
profil et une logique d'affichage impossible à tenir. Ce document pose donc d'abord la
fondation commune, et rien d'autre.

\begin{note}[Comment lire ce document]
Il est organisé en deux parties. La \textbf{partie I} décrit la fonctionnalité telle que
l'utilisateur et le client la voient, sans vocabulaire technique. La \textbf{partie II}
décrit comment on la construit. Les deux couvrent le même périmètre ; on peut valider la
partie I sans lire la partie II.
\end{note}

\begin{constat}[Méthode]
Toutes les affirmations sur l'existant de ce document ont été vérifiées en lisant le code
des trois dépôts. Les chemins de fichiers cités sont réels. Là où une information manque,
c'est écrit.
\end{constat}

\tableofcontents

% =====================================================================
\part{Approche naïve}
% =====================================================================

% Fragment partage : inclus par socle-compte.tex (dossier autonome)
% et par dossier-conception-alanya.tex (dossier client assemble).
% Ne pas y mettre de preambule ni de \part : c'est l'appelant qui les pose.

\chapter{Ce que l'on ajoute, et pourquoi}

\section{La situation aujourd'hui}

Dans Alanya, tous les comptes sont identiques. Un commerçant, une personnalité publique et
un utilisateur ordinaire ont exactement le même profil et la même apparence dans une liste
de conversations. Rien ne permet de distinguer l'entreprise officielle d'un imitateur qui
aurait pris le même nom et la même photo.

C'est ce manque que les trois fonctionnalités demandées viennent combler --- et c'est
pourquoi elles partagent la même fondation.

\section{Deux questions, pas une}

L'erreur naturelle serait de créer une liste de « genres de comptes » : personnel,
business, vérifié, officiel, business vérifié\ldots{} Cette liste devient ingérable dès
qu'on ajoute un cas.

On sépare donc deux questions \textbf{indépendantes} :

\begin{center}
\begin{tabularx}{\linewidth}{@{}p{4.2cm}X@{}}
\toprule
\thd{Quel genre de compte ?} & Personnel, business, ou officiel. C'est ce que le compte
\emph{est}. \\
\thd{A-t-il été vérifié ?} & Aucune démarche, en attente, vérifié, refusé, révoqué, ou
expiré. C'est ce qu'Alanya \emph{atteste} à son sujet. \\
\bottomrule
\end{tabularx}
\end{center}

Ces deux questions ne dépendent pas l'une de l'autre. Un commerce peut exister sans être
vérifié. Une personne privée peut être vérifiée sans être un commerce --- c'est le cas
d'une personnalité publique. Le compte officiel Alanya, lui, est la combinaison
« officiel + vérifié d'office ».

\begin{note}[Le bénéfice concret]
Poser deux questions séparées plutôt qu'une seule liste, c'est ce qui permet d'ajouter
demain un nouveau genre de compte, ou un nouvel état de vérification, sans toucher à
l'autre axe ni refaire les écrans.
\end{note}

\section{Ce que couvre l'étape 1 --- et ce qu'elle ne couvre pas}

\begin{tabularx}{\linewidth}{@{}p{5.4cm}X@{}}
\toprule
\thd{Dans cette étape} & Savoir \emph{qui} est un compte : son genre, son état de
vérification, et le badge qui en découle partout dans l'application. \\
\thd{Étapes suivantes} & Ce que ce compte \emph{contient} : la fiche business (horaires,
adresse, catalogue, liens), le dossier de certification (pièces justificatives,
approbation, abonnement), et la diffusion du compte officiel. \\
\bottomrule
\end{tabularx}

La raison de cette coupure est simple : le genre et l'état d'un compte suffisent déjà à
afficher le bon badge partout. Le contenu détaillé d'une fiche business dépend de règles
qu'on n'a pas encore écrites --- le figer maintenant obligerait à tout refaire à l'étape
suivante.

\chapter{Ce que l'utilisateur voit}

\section{Le principe du badge}

Un badge qui combinerait « c'est un commerce » et « c'est vérifié » en un seul dessin
devient illisible à la taille où il apparaît réellement --- environ 12 pixels, à côté d'un
nom dans une liste de conversations. Et deux badges côte à côte mangent la largeur du nom,
qui est déjà tronqué sur un petit écran.

La solution retenue sépare les deux informations sur deux canaux visuels distincts :

\begin{center}
\begin{tabularx}{0.92\linewidth}{@{}p{4.6cm}X@{}}
\toprule
\thd{Le glyphe dit \emph{quoi}} & Une coche pour une personne notable, un panier pour un
commerce, la marque Alanya pour le compte officiel. \\
\thd{Le contenant dit \emph{vérifié ou pas}} & Un glyphe nu = déclaré par le compte
lui-même. Un glyphe posé dans un \textbf{sceau festonné} \badgeSceauNu[0.9] = attesté par
Alanya. \\
\bottomrule
\end{tabularx}
\end{center}

\subsection*{Pourquoi un sceau festonné, et pas un rond}

Un disque plein ne se distingue d'aucune autre icône circulaire de l'interface : avatars,
puces de comptage, indicateurs de présence. À 12~pixels, un badge rond devient un point
coloré parmi d'autres points colorés.

Le feston à douze lobes, lui, a une silhouette reconnaissable même quand le glyphe qu'il
contient n'est plus lisible. C'est la propriété qui compte : l'utilisateur qui distingue
la \emph{forme} sait qu'Alanya s'est prononcé sur ce compte, sans avoir à lire le
détail. Et cette forme est bien plus difficile à imiter avec un caractère Unicode dans un
nom d'affichage qu'un simple rond ou une simple coche.

\section{Les cinq rendus}

\begin{center}
\setlength{\extrarowheight}{7pt}
\begin{tabular}{@{}llc@{\hspace{1.6cm}}c@{}}
\toprule
\thd{Genre} & \thd{État} & \thd{À taille réelle} & \thd{Agrandi} \\
\midrule
Personnel & non vérifié & --- & --- \\
Personnel & vérifié     & \badgePastilleCoche[1]  & \badgePastilleCoche[2.2] \\
Business  & déclaré     & \badgePanier[1]         & \badgePanier[2.2] \\
Business  & vérifié     & \badgePastillePanier[1] & \badgePastillePanier[2.2] \\
Officiel  & ---         & \badgeMarque[1]         & \badgeMarque[2.2] \\
\bottomrule
\end{tabular}
\end{center}

\medskip

La colonne « à taille réelle » est là pour être jugée telle quelle : c'est exactement la
taille à laquelle le badge apparaîtra à côté d'un nom. Si un rendu n'est pas
distinguable à cette échelle, il doit être redessiné avant l'implémentation.

\begin{note}[Pourquoi une couleur différente pour l'officiel]
Le sceau du compte officiel est doré, pas indigo. Il n'existe qu'un seul compte
officiel dans toute l'application : lui donner une couleur à part évite qu'il soit confondu
avec un compte simplement vérifié, et rend l'usurpation visuellement plus difficile.
\end{note}

\section{Où le badge apparaît}

Le même badge, dessiné par le même composant, dans quatre contextes :

\begin{itemize}[leftmargin=1.4em]
  \item la \textbf{liste de conversations}, juste après le nom ;
  \item l'\textbf{en-tête d'une conversation} ouverte ;
  \item la \textbf{page de profil} du compte ;
  \item les \textbf{résultats de recherche} et les listes de contacts.
\end{itemize}

\begin{center}
  \imageOuCadre{maquettes/liste-conversations-badges.png}{0.86\linewidth}%
    {Liste de conversations montrant les cinq variantes de badge en situation\\
     \texttt{maquettes/liste-conversations-badges.png}}
\end{center}

\begin{note}[Maquette interactive]
La version vivante de cette maquette permet de faire varier la taille du badge de 10 à
28~pixels pour juger la lisibilité, et montre la conversation officielle avec son bandeau
de lecture seule :\\
\href{https://claude.ai/code/artifact/029627c6-5dae-4b2a-a86b-3e28c899d824}%
{\texttt{claude.ai/code/artifact/029627c6}}\\
Source versionnée : \code{docs/architecture/maquettes/liste-conversations-badges.html}
\end{note}

\section{Pourquoi la distinction déclaré / vérifié est critique}

Se déclarer « commerce » prend trente secondes et n'engage personne : le compte l'affirme
lui-même, personne ne l'a contrôlé. Être vérifié suppose au contraire un dossier examiné
par Alanya.

Si les deux badges avaient le même poids visuel, n'importe qui créerait un compte business
et hériterait d'une caution qui ne lui a jamais été accordée. C'est exactement le scénario
qu'un arnaqueur cherche.

\begin{alerte}[La règle qui en découle]
Le sceau festonné est la signature d'Alanya. Rien d'autre dans l'interface ne doit pouvoir lui
ressembler --- ni une icône décorative, ni un emoji dans un nom d'utilisateur.
\end{alerte}

\section{Protéger le badge des contournements}

Un système de badge peut être contourné sans jamais être attaqué techniquement : il suffit
qu'un utilisateur s'appelle « Alanya~\checkmark{} » pour que le nom affiché imite la
certification.

Deux protections sont donc posées dès cette étape :

\begin{enumerate}[leftmargin=1.6em]
  \item les caractères de type coche, étoile et sceau sont \textbf{refusés} dans le nom et
        le pseudo ;
  \item les variantes du nom « Alanya » sont \textbf{réservées} et ne peuvent pas être
        prises par un compte ordinaire.
\end{enumerate}

\chapter{Les règles de vie d'un compte}

\section{Devenir un compte business, et revenir en arrière}

Le passage d'un compte personnel à un compte business est \textbf{réversible}. Les règles
proposées :

\begin{center}
\begin{tabularx}{\linewidth}{@{}p{5.2cm}X@{}}
\toprule
\thd{Personnel $\rightarrow$ business} & Libre, à l'initiative de l'utilisateur. Devenir
commerce est déclaratif : cela n'atteste de rien et ne demande donc aucune validation. \\
\thd{Business $\rightarrow$ personnel} & Libre également. La fiche business est
\textbf{conservée mais masquée}, pas supprimée : un retour en arrière ne doit pas détruire
un catalogue de produits saisi pendant des semaines. \\
\thd{Une vérification en cours} & La demande est annulée : elle portait sur une identité
d'entreprise qui n'est plus revendiquée. \\
\thd{Une vérification obtenue} & Elle est \textbf{révoquée}, pas suspendue. Le badge
atteste d'une entreprise ; si le compte n'en est plus une, l'attestation n'a plus d'objet.
Un nouveau dossier sera nécessaire pour la retrouver. \\
\bottomrule
\end{tabularx}
\end{center}

\begin{note}[À confirmer avec le client]
La bascule business $\rightarrow$ personnel pourrait aussi être soumise à validation d'un
administrateur, pour éviter qu'un commerce échappe à un signalement en repassant en
personnel. La règle proposée ci-dessus est la plus simple ; l'alternative est plus sûre.
\end{note}

\section{La vie d'une vérification}

\begin{center}
\begin{tikzpicture}[
  x=1cm, y=1cm,
  etat/.style={draw=AlanyaLine, fill=AlanyaSurface, rounded corners=3pt,
               font=\sffamily\small, inner sep=6pt, minimum height=0.8cm,
               text=AlanyaInk},
  etatOk/.style={etat, draw=AlanyaSuccess, line width=0.9pt, fill=AlanyaSuccess!8},
  etatKo/.style={etat, draw=danger, fill=danger!6},
  fleche/.style={-{Stealth[length=2.2mm]}, draw=AlanyaGray, semithick},
  etiq/.style={font=\sffamily\scriptsize, text=AlanyaGray,
               fill=white, inner sep=1.5pt}
]
  \node[etat]   (aucun)   at (0,0)      {Aucune démarche};
  \node[etat]   (attente) at (4.3,0)    {En attente};
  \node[etatOk] (ok)      at (8.3,0)    {Vérifié};
  \node[etatKo] (revoque) at (11.8,0)   {Révoqué};
  \node[etatKo] (refuse)  at (1.7,-2.1) {Refusé};
  \node[etatKo] (expire)  at (8.3,-2.1) {Expiré};

  % Chemin nominal
  \draw[fleche] (aucun)   -- node[etiq, above] {dépôt}       (attente);
  \draw[fleche] (attente) -- node[etiq, above] {approbation} (ok);
  \draw[fleche] (ok)      -- node[etiq, above] {abus}        (revoque);
  \draw[fleche] (ok)      -- node[etiq, right] {échéance}    (expire);

  % Rejet
  \draw[fleche] (attente.south west) -- node[etiq, pos=0.55, above left, inner sep=1pt]
    {rejet} (refuse.north east);

  % Retours : deux couloirs qui ne se croisent pas
  \draw[fleche] (refuse.north west) -- node[etiq, pos=0.5, above left, inner sep=1pt]
    {nouveau dossier} (aucun.south);

  \draw[fleche] (expire.west) -- (4.3,-2.1)
    node[etiq, pos=0.55, below] {renouvellement} -- (attente.south);
\end{tikzpicture}
\end{center}

\section{L'expiration et sa relance}

Une vérification n'est pas éternelle : elle a une date de fin. Le comportement attendu :

\begin{center}
\begin{tabularx}{\linewidth}{@{}p{3.4cm}X@{}}
\toprule
\thd{À J\,--\,30} & Le titulaire est prévenu qu'il doit renouveler, par trois canaux à la
fois : une notification sur son téléphone, un e-mail, et un avertissement visible dans
l'application. L'avertissement n'est envoyé qu'une fois. \\
\thd{À l'échéance} & La vérification est \textbf{suspendue}. Le badge disparaît. Le compte
et son contenu restent intacts --- seule l'attestation tombe. \\
\thd{Après} & Le titulaire peut déposer un nouveau dossier pour retrouver son badge. \\
\bottomrule
\end{tabularx}
\end{center}

\begin{note}[Suspendu, pas supprimé]
Une vérification expirée est un état distinct d'un refus. Le compte a bien été vérifié par
le passé ; il a simplement laissé passer la date. Cette nuance compte pour le service
client, et pour ne pas retraiter un dossier déjà instruit.
\end{note}

\section{Le compte officiel}

Le compte officiel est le compte au nom de l'application. Ses règles :

\begin{itemize}[leftmargin=1.4em]
  \item il est \textbf{créé par un administrateur}, jamais par une inscription ordinaire ;
  \item sa photo de profil est le \textbf{logo Alanya} ;
  \item il porte le badge officiel, distinct du badge de vérification ;
  \item les utilisateurs \textbf{ne peuvent pas lui répondre}. Ses messages sont des
        diffusions, pas des conversations.
\end{itemize}

Concrètement, dans une conversation avec le compte officiel, la zone de saisie est
remplacée par un bandeau explicatif, et les actions « répondre » et « réagir » n'apparaissent
pas sur ses messages.

\begin{alerte}[Ce n'est pas qu'une question d'affichage]
Masquer la zone de saisie ne suffit pas : un client modifié pourrait envoyer le message
quand même. Le refus doit être appliqué par le serveur. La partie II décrit où.
\end{alerte}

\section{Un compte business ne peut pas être administrateur}

Décision prise : les rôles d'administration de la plateforme et les genres de compte sont
mutuellement exclusifs. Un compte business ou officiel ne peut pas recevoir de droits
d'administration, et réciproquement.

La raison est de sécurité : un compte visible du public, susceptible d'être usurpé ou
compromis, ne doit pas ouvrir la porte au back-office.


% =====================================================================
\part{Approche technique}
% =====================================================================

% Fragment partage : inclus par socle-compte.tex (dossier autonome)
% et par dossier-conception-alanya.tex (dossier client assemble).
% Ne pas y mettre de preambule ni de \part : c'est l'appelant qui les pose.

\chapter{Modèle de données}

\section{Pourquoi étendre \code{users} plutôt que créer une table \code{accounts}}

La conception initiale prévoyait une nouvelle table \code{accounts} servant de socle. La
lecture du schéma réel invalide cette approche.

\begin{constat}
\textbf{14 tables} référencent \code{users(alanyaID)} par clé étrangère :
\code{preferredContact}, \code{blocked}, \code{userAccess}, \code{conversation},
\code{conv\_participants}, \code{message}, \code{statut}, \code{statut\_views},
\code{meeting}, \code{participant}, \code{callHistory}, \code{reserved\_alanya\_phone},
\code{user\_push\_devices} et \code{user\_notification\_prefs} --- soit 34 contraintes au
total (sources : \code{migrations/001\_initial\_schema.sql}, \code{010}, \code{020},
\code{021}).
\end{constat}

Introduire une table parallèle imposerait de réécrire ces 34 contraintes. On étend donc
\code{users}, qui \emph{est} le socle de compte.

Ce choix a un second bénéfice, qui n'est pas un hasard : la conception exige que le genre
et l'état de vérification soient lisibles dans la requête de liste de conversations
\textbf{sans jointure supplémentaire}. Cette requête joint déjà \code{users} pour récupérer
le nom et l'avatar du correspondant
(\code{src/controllers/conversationController.js}, ligne~24). Deux colonnes de plus sur
\code{users} sont donc gratuites ; une table \code{accounts} aurait imposé une jointure de
plus sur le chemin le plus chaud de l'application.

\section{La migration}

\begin{lstlisting}[language=SQL, caption={\code{migrations/023\_account\_socle.sql}}]
-- 023_account_socle.sql -- socle de compte unifie
--
-- Deux axes orthogonaux portes par `users` :
--   account_type        : ce que le compte EST
--   verification_status : ce qu'Alanya ATTESTE a son sujet
--
-- TINYINT et non ENUM : ajouter une valeur a un ENUM impose un ALTER TABLE
-- sur une table qui grossit. L'enumeration vit cote Node.

ALTER TABLE users
  ADD COLUMN account_type TINYINT NOT NULL DEFAULT 0
    COMMENT '0=personnel 1=business 2=officiel',
  ADD COLUMN verification_status TINYINT NOT NULL DEFAULT 0
    COMMENT '0=aucun 1=en_attente 2=verifie 3=refuse 4=revoque 5=expire',
  ADD COLUMN verified_until DATETIME NULL
    COMMENT 'Echeance de la verification (NULL si non verifie)',
  ADD COLUMN verification_reminder_sent TINYINT NOT NULL DEFAULT 0
    COMMENT '1 = relance J-30 deja envoyee pour l echeance courante';

-- Annuaire business a venir : filtrer par genre puis par etat.
ALTER TABLE users
  ADD INDEX idx_users_account_type (account_type, verification_status);

-- Balayage du scheduler d'expiration : ne remonter que les comptes verifies
-- ayant une echeance, sans parcourir toute la table.
ALTER TABLE users
  ADD INDEX idx_users_verified_until (verified_until);
\end{lstlisting}

\begin{note}[Aucune colonne \code{badge\_type}]
Le badge n'est jamais stocké. Il se calcule à l'affichage à partir des deux axes. Une
valeur stockée finirait désynchronisée le jour où une certification expire ou est révoquée
--- ce qui est précisément le cas que le système doit gérer.
\end{note}

\section{Les énumérations, côté Node}

\begin{lstlisting}[language=JavaScript, caption={\code{src/constants/accountTypes.js} (nouveau)}]
const ACCOUNT_TYPE = { PERSONNEL: 0, BUSINESS: 1, OFFICIEL: 2 };

const VERIFICATION = {
  AUCUN:      0,
  EN_ATTENTE: 1,
  VERIFIE:    2,
  REFUSE:     3,
  REVOQUE:    4,
  EXPIRE:     5,
};

module.exports = { ACCOUNT_TYPE, VERIFICATION };
\end{lstlisting}

Le dossier \code{src/constants/} existe déjà (\code{participantLimits.js}) : la convention
est respectée.

\section{Raccordement au schéma existant}

\begin{center}
\begin{tikzpicture}[
  tbl/.style={draw=AlanyaLine, fill=white, rounded corners=3pt,
              font=\sffamily\scriptsize, inner sep=5pt, text=AlanyaInk},
  socle/.style={draw=AlanyaIndigo, line width=1pt, fill=AlanyaContainer,
                rounded corners=4pt, font=\sffamily\small\bfseries,
                inner sep=8pt, text=AlanyaIndigoDark, align=center},
  futur/.style={draw=AlanyaMute, dashed, fill=white, rounded corners=3pt,
                font=\sffamily\scriptsize\itshape, inner sep=5pt, text=AlanyaGray},
  lien/.style={draw=AlanyaLine, thick}
]
  \node[socle] (users) at (0,0) {users\\[2pt]
    \footnotesize\mdseries account\_type\\
    \footnotesize\mdseries verification\_status\\
    \footnotesize\mdseries verified\_until};

  \node[tbl] (a) at (-4.6,2.1) {message};
  \node[tbl] (b) at (-4.6,1.1) {conversation};
  \node[tbl] (c) at (-4.6,0.1) {conv\_participants};
  \node[tbl] (d) at (-4.6,-0.9) {preferredContact};
  \node[tbl] (e) at (-4.6,-1.9) {statut};
  \node[tbl] (f) at (-4.6,-2.9) {\ldots{} 9 autres};

  \node[futur] (g) at (4.7,1.5)  {business\_profiles};
  \node[futur] (h) at (4.7,0.4)  {verification\_requests};
  \node[futur] (i) at (4.7,-0.7) {contact\_list};

  \foreach \n in {a,b,c,d,e,f} { \draw[lien] (\n.east) -- (users.west); }
  \foreach \n in {g,h,i} { \draw[lien, dashed, draw=AlanyaMute] (users.east) -- (\n.west); }

  \node[font=\sffamily\scriptsize, text=AlanyaGray] at (-4.6,3.0) {14 tables existantes};
  \node[font=\sffamily\scriptsize, text=AlanyaGray] at (4.7,2.4) {étapes suivantes};
\end{tikzpicture}
\end{center}

\chapter{Dérivation du badge}

\section{Une fonction pure, trois implémentations identiques}

Le badge est une fonction de deux entiers vers une variante d'affichage. Aucun accès base,
aucune requête réseau.

\begin{lstlisting}[language=Dart, caption={\code{lib/widgets/common/account\_badge.dart} (nouveau) --- forme cible}]
enum AccountBadge { aucun, cocheVerifiee, panierDeclare, panierVerifie, officiel }

AccountBadge resolveAccountBadge(int accountType, int verificationStatus) {
  final verifie = verificationStatus == 2; // VERIFIE
  switch (accountType) {
    case 2: return AccountBadge.officiel;
    case 1: return verifie ? AccountBadge.panierVerifie : AccountBadge.panierDeclare;
    default: return verifie ? AccountBadge.cocheVerifiee : AccountBadge.aucun;
  }
}
\end{lstlisting}

\begin{center}
\begin{tabular}{@{}llll@{}}
\toprule
\thd{\code{account\_type}} & \thd{\code{verification\_status}} & \thd{Variante} & \thd{Rendu} \\
\midrule
0 & $\neq 2$ & \code{aucun}          & --- \\
0 & $= 2$    & \code{cocheVerifiee}  & \badgePastilleCoche[1] \\
1 & $\neq 2$ & \code{panierDeclare}  & \badgePanier[1] \\
1 & $= 2$    & \code{panierVerifie}  & \badgePastillePanier[1] \\
2 & indifférent & \code{officiel}    & \badgeMarque[1] \\
\bottomrule
\end{tabular}
\end{center}

\begin{alerte}[Un seul widget, pas trois]
Le composant doit être unique et utilisé aux quatre endroits listés en partie~I. Trois
implémentations concurrentes divergeront en deux semaines --- c'est exactement ce qui est
arrivé aux quatre indigos concurrents que le design system a dû unifier
(\code{lib/core/theme/app\_colors.dart}).
\end{alerte}

\section{Le piège à corriger avant tout}

\begin{constat}[Vérifié dans le code]
Le champ \code{typeCompte} du modèle \code{User} est \textbf{codé en dur à \code{0}} partout
où un \code{User} est reconstruit depuis une charge utile partielle :
\code{lib/core/services/call\_service.dart} ligne~207,
\code{lib/widgets/chat/contact\_message\_preview.dart} ligne~46, et
\code{lib/screens/chats/contact\_detail\_screen.dart} ligne~104.
\end{constat}

Les nouveaux champs suivront le même chemin et subiront le même sort : un badge dérivé
afficherait « compte personnel non vérifié » à ces trois endroits, quel que soit le compte
réel. Ces trois sites doivent être corrigés dans le même lot que l'ajout du badge, sinon le
défaut est invisible en développement et visible en production.

\section{Composant existant à ne pas confondre}

\code{lib/widgets/common/app\_badge.dart} existe déjà mais n'a aucun rapport : c'est
\code{CountBadge}, la pastille de comptage des messages non lus. Le nouveau composant est
donc un fichier distinct, \code{account\_badge.dart}.

\chapter{Règles appliquées par le serveur}

\section{Exclusion business / administrateur}

Règle : \code{account\_type != 0} $\Rightarrow$ \code{type\_compte = 0}.

\begin{note}[Pourquoi deux colonnes distinctes]
\code{users.type\_compte} existe déjà et porte le \textbf{rôle d'administration}
(\code{0} utilisateur, \code{>= 1} admin, \code{>= 2} super-admin). Il est lu par
\code{src/middleware/adminAuth.js} (lignes~41 et~60),
\code{lib/providers/admin\_provider.dart} (lignes~245--246) et
\code{talky-admin/lib/auth.ts}. Réutiliser cette colonne pour le genre de compte aurait
cassé l'administration existante ; d'où la colonne séparée \code{account\_type}.
\end{note}

Le contrôle est à poser aux deux endroits qui écrivent ces champs, dans
\code{src/controllers/admin/users.js} : \code{setAccountType} (rôle admin, existant) et le
nouveau point d'entrée qui fixera \code{account\_type}. Rejet en 409 avec un message
explicite plutôt qu'un écrasement silencieux.

\section{Validation des noms d'affichage}

Le filtrage doit vivre \textbf{côté serveur}, et non dans l'application mobile : le mobile
n'est pas le seul écrivain. L'administration crée aussi des comptes
(\code{src/controllers/admin/userCreate.js}), et la modification de profil passe par
l'API.

\begin{lstlisting}[language=JavaScript, caption={\code{src/utils/displayNameGuard.js} (nouveau) --- principe}]
// Blocs Unicode a refuser dans `nom` et `pseudo` :
//   U+2705, U+2713-U+2714 (coches)        U+2B50, U+2606, U+2B51 (etoiles)
//   U+1F396-U+1F397, U+1F3C5 (medailles)  U+269C (fleur de lys / sceau)
//   Variation selectors utilises pour composer un faux badge.
//
// Plus : refus des variantes normalisees du nom "Alanya", en tenant compte
//   de la casse, des accents, des espaces inseres, et des homoglyphes
//   cyrilliques U+0410 U+0430 (a), U+0435 (e), U+043E (o), U+0443 (y),
//   U+0441 (c) -- visuellement identiques, distincts pour la base.
\end{lstlisting}

\begin{alerte}[Les homoglyphes]
Refuser la chaîne « Alanya » ne suffit pas. Le même mot dont le \emph{A} initial est un A
cyrillique (U+0410) au lieu du A latin (U+0041) est une chaîne \emph{différente} pour la
base de données, et rigoureusement identique à l'{\oe}il. La comparaison doit donc se faire
sur une forme normalisée, après repli des homoglyphes.
\end{alerte}

\section{Interdire la réponse au compte officiel}

\begin{constat}[Un seul point d'ancrage suffit]
Les deux chemins d'envoi de message --- REST
(\code{src/controllers/messageController.js} ligne~178) et Socket.IO
(\code{src/socket/handlers/chat/messageSend.js} ligne~142) --- appellent tous deux la même
fonction \code{evaluateDirectMessageSend} de \code{src/utils/blockUtils.js}. Une seule
modification couvre donc les deux.
\end{constat}

Cette fonction résout déjà le correspondant d'une conversation 1-1 et renvoie une action
(\code{deliver}, \code{reject} ou \code{silent}). Il suffit d'y ajouter un cas :

\begin{lstlisting}[language=JavaScript, caption={\code{src/utils/blockUtils.js} --- ajout dans \code{evaluateDirectMessageSend}}]
const evaluateDirectMessageSend = async (conversationID, senderID) => {
  const peerId = await getDirectConversationPeer(conversationID, senderID);
  if (peerId == null) return { isDirect: false };

  // Nouveau : le compte officiel diffuse, il ne converse pas.
  const [rows] = await pool.execute(
    'SELECT account_type FROM users WHERE alanyaID = ?', [peerId],
  );
  if (rows[0] && rows[0].account_type === ACCOUNT_TYPE.OFFICIEL) {
    return { isDirect: true, peerId, action: 'reject', code: 'OFFICIAL_READONLY' };
  }

  // ... suite inchangee : evaluation du blocage
};
\end{lstlisting}

Côté mobile, la zone de saisie de \code{lib/screens/chats/chat/chat\_input.dart} est
remplacée par un bandeau, et les actions « répondre » / « réagir » sont retirées des bulles.

\begin{note}[À trancher à l'étape 3]
Faut-il une conversation en lecture seule, ou un canal de diffusion distinct des
conversations ? Le choix dépend de la stratégie de diffusion, qui est le sujet de l'étape
sur le compte officiel. La règle serveur ci-dessus vaut dans les deux cas.
\end{note}

\section{Expiration : réutiliser le mécanisme existant}

Le comportement demandé --- relance à J\,--\,30 puis suspension --- a déjà son patron exact
dans le code.

\begin{center}
\begin{tabularx}{\linewidth}{@{}p{5.1cm}X@{}}
\toprule
\thd{Ordonnanceur} & \code{src/services/meetingScheduler.js} : un \code{setInterval} d'une
minute, une requête qui remonte les échéances proches, et une colonne \code{reminder\_sent}
pour ne jamais notifier deux fois. Le squelette se transpose tel quel en
\code{verificationScheduler.js}. \\
\thd{E-mail} & \code{src/services/mailService.js} : \code{sendMail} et
\code{renderHtmlEmail} sur le gabarit \code{src/templates/email-template.html}, déjà en
service pour l'OTP de réinitialisation. \\
\thd{Notification} & \code{src/services/notificationService.js}, déjà utilisé par
l'ordonnanceur de réunions. \\
\thd{Avertissement in-app} & Dérivé de \code{verified\_until}, déjà présent dans la charge
utile utilisateur --- aucun appel supplémentaire. \\
\bottomrule
\end{tabularx}
\end{center}

\begin{lstlisting}[language=SQL, caption={Les deux requêtes du scheduler}]
-- 1. Relance J-30, une seule fois par echeance
SELECT alanyaID, nom, email, verified_until
FROM users
WHERE verification_status = 2
  AND verified_until IS NOT NULL
  AND verified_until > UTC_TIMESTAMP()
  AND verified_until <= DATE_ADD(UTC_TIMESTAMP(), INTERVAL 30 DAY)
  AND verification_reminder_sent = 0;

-- 2. Suspension a l'echeance
UPDATE users
SET verification_status = 5,            -- EXPIRE
    verification_reminder_sent = 0      -- rearme pour un futur renouvellement
WHERE verification_status = 2
  AND verified_until IS NOT NULL
  AND verified_until <= UTC_TIMESTAMP();
\end{lstlisting}

\begin{note}[Pourquoi pas de file de jobs ici]
La conception initiale demandait une file de jobs (BullMQ / Redis) dès le départ. Elle
n'existe pas aujourd'hui dans le backend. Pour les relances de vérification, le volume est
faible et espacé : un ordonnanceur en processus suffit, et il a déjà fait ses preuves sur
les rappels de réunion. La file redeviendra nécessaire au fan-out des diffusions du compte
officiel --- c'est là qu'il faudra la traiter, pas ici.
\end{note}

\chapter{Impact fichier par fichier}

\section{Backend --- Alanya-Backend}

\begin{center}
\begin{tabularx}{\linewidth}{@{}p{6.4cm}p{1.6cm}X@{}}
\toprule
\thd{Fichier} & \thd{Nature} & \thd{Objet} \\
\midrule
\code{migrations/023\_account\_socle.sql} & nouveau & Les quatre colonnes et deux index \\
\code{src/constants/accountTypes.js} & nouveau & Énumérations partagées \\
\code{src/utils/displayNameGuard.js} & nouveau & Filtrage des noms, réservation d'« Alanya » \\
\code{src/services/verificationScheduler.js} & nouveau & Relance J\,--\,30 et suspension \\
\code{src/utils/blockUtils.js} & modifié & Refus de réponse au compte officiel \\
\code{src/controllers/userController.js} & modifié & Exposer les deux axes ; appliquer le filtrage des noms \\
\code{src/controllers/conversationController.js} & modifié & Ajouter les deux axes à la projection du correspondant \\
\code{src/controllers/preferredContactController.js} & modifié & Idem sur la liste de contacts \\
\code{src/controllers/messageController.js} & modifié & Idem sur l'expéditeur d'un message \\
\code{src/controllers/admin/users.js} & modifié & Piloter \code{account\_type} ; exclusion business/admin ; filtre par genre \\
\code{src/controllers/admin/userCreate.js} & modifié & Création du compte officiel ; filtrage des noms \\
\code{server.js} & modifié & Démarrage du \code{verificationScheduler} \\
\bottomrule
\end{tabularx}
\end{center}

\section{Mobile --- Talky (Flutter)}

\begin{center}
\begin{tabularx}{\linewidth}{@{}p{6.4cm}p{1.6cm}X@{}}
\toprule
\thd{Fichier} & \thd{Nature} & \thd{Objet} \\
\midrule
\code{lib/widgets/common/account\_badge.dart} & nouveau & Le composant unique et sa dérivation \\
\code{lib/talky\_models.dart} & modifié & \code{accountType}, \code{verificationStatus}, \code{verifiedUntil} sur \code{User} \\
\code{lib/core/db/app\_database.dart} & modifié & Colonnes sur \code{LocalUsers}, \code{schemaVersion} et \code{onUpgrade} \\
\code{lib/core/services/local\_cache\_repository.dart} & modifié & Propager les champs dans le cache \\
\code{lib/core/services/call\_service.dart} & \textbf{corrigé} & \code{typeCompte: 0} codé en dur, ligne~207 \\
\code{lib/widgets/chat/contact\_message\_preview.dart} & \textbf{corrigé} & Idem, ligne~46 \\
\code{lib/screens/chats/contact\_detail\_screen.dart} & \textbf{corrigé} & Idem, ligne~104 \\
\code{lib/screens/chats/chats\_screen.dart} & modifié & Badge dans la liste de conversations \\
\code{lib/screens/chats/chat\_detail\_screen.dart} & modifié & Badge dans l'en-tête ; masquage de la saisie \\
\code{lib/screens/chats/chat/chat\_input.dart} & modifié & Bandeau « diffusion » à la place du champ \\
\code{lib/screens/chats/chat/chat\_bubbles.dart} & modifié & Retrait de « répondre » et « réagir » \\
\code{lib/l10n/*} & modifié & Libellés des badges, du bandeau et des avertissements \\
\bottomrule
\end{tabularx}
\end{center}

\begin{alerte}[Régénération obligatoire]
Toute modification de \code{app\_database.dart} impose
\code{dart run build\_runner build -{}-delete-conflicting-outputs}, et la montée de
\code{schemaVersion} doit passer par \code{onUpgrade} --- pas par une recréation, sous peine
de vider le cache local des utilisateurs déjà installés.
\end{alerte}

\section{Administration --- Alanya Admin}

\begin{center}
\begin{tabularx}{\linewidth}{@{}p{6.4cm}p{1.6cm}X@{}}
\toprule
\thd{Fichier} & \thd{Nature} & \thd{Objet} \\
\midrule
\code{types/index.ts} & modifié & \code{accountType}, \code{verificationStatus}, \code{verifiedUntil} \\
\code{app/(dashboard)/users/page.tsx} & modifié & Colonne badge et filtre par genre de compte \\
\code{app/(dashboard)/users/[id]/page.tsx} & modifié & Pilotage du genre et de l'état \\
\code{app/(dashboard)/users/new/page.tsx} & modifié & Création du compte officiel \\
\code{components/ui/badge.tsx} & modifié & Variantes de badge de compte \\
\bottomrule
\end{tabularx}
\end{center}

\begin{note}[Conversion des noms]
Le client HTTP de l'administration convertit automatiquement les clés
\code{snake\_case} en \code{camelCase} (\code{lib/api.ts}). Les colonnes
\code{account\_type} et \code{verification\_status} arriveront donc en
\code{accountType} et \code{verificationStatus} sans travail supplémentaire.
\end{note}

\chapter{Points de vigilance et questions ouvertes}

\section{À vérifier avant d'appliquer la migration}

\begin{alerte}[Le backend n'a pas d'ordonnanceur de migrations]
Il n'existe ni exécuteur, ni table de suivi : les fichiers de \code{migrations/} sont
appliqués \textbf{à la main}. Rien ne garantit que l'état réel de la base corresponde au
contenu du dossier. Avant d'appliquer 023, il faut constater l'état réel des colonnes.
\end{alerte}

\begin{alerte}[La migration 013 contient du SQL invalide]
\code{migrations/013\_view\_once.sql} comporte un point-virgule au milieu de l'instruction
\code{ALTER TABLE} : la seconde colonne, \code{viewedAt}, n'a jamais pu s'appliquer par ce
fichier. Elle figure pourtant dans le schéma consolidé
\code{docs/schema\_complet\_alanyBD2027.sql}. Il faut trancher lequel des deux décrit la
production.
\end{alerte}

\section{Collision de numérotation}

Le numéro 023 est également revendiqué par \code{docs/listes-contacts-plan.md}, qui prévoit
\code{023\_contact\_lists.sql}. Le socle prenant 023, les listes de contacts passent à
\textbf{024}. Le document doit être corrigé avant sa conversion, sinon les deux migrations
entreront en collision à l'implémentation.

\section{Écart assumé avec la conception initiale}

\begin{center}
\begin{tabularx}{\linewidth}{@{}p{4.6cm}X@{}}
\toprule
\thd{Conception initiale} & \thd{Ce document} \\
\midrule
Table \code{accounts} dédiée & Extension de \code{users} --- 14 tables et 34 contraintes
en dépendent déjà \\
« Fuir le type \code{ENUM} » & Appliqué : \code{TINYINT} partout. À noter que le schéma
existant en contient déjà (\code{migrations/020} et \code{021}) ; on ne revient pas dessus \\
File de jobs dès la conception & Reportée : l'ordonnanceur en processus suffit pour les
relances. La file redevient nécessaire au fan-out des diffusions \\
\bottomrule
\end{tabularx}
\end{center}

\section{Ce qui reste à décider}

\begin{enumerate}[leftmargin=1.6em]
  \item La bascule business $\rightarrow$ personnel est-elle libre, ou soumise à validation
        d'un administrateur ?
  \item Un compte business dont la vérification a expiré reste-t-il visible dans l'annuaire
        business, sans badge ?
  \item La durée par défaut d'une vérification --- un an est l'usage, mais ce n'est pas
        tranché.
  \item Y a-t-il un seul compte officiel, ou plusieurs (support, actualités\ldots) ? Le
        modèle proposé en supporte plusieurs sans modification.
\end{enumerate}


\vfill
\begin{center}\small\color{AlanyaGray}\sffamily
Fin du dossier --- étape 1. \\
Étape suivante : listes de contacts préférés.
\end{center}

\end{document}
